
\documentclass{article}
\usepackage{colortbl}
\usepackage{makecell}
\usepackage{multirow}
\usepackage{supertabular}

\begin{document}

\newcounter{utterance}

\centering \large Interaction Transcript for game `hot\_air\_balloon', experiment `air\_balloon\_survival\_de\_complexity\_easy', episode 0 with gpt{-}oss{-}120b{-}t1.0.
\vspace{24pt}

{ \footnotesize  \setcounter{utterance}{1}
\setlength{\tabcolsep}{0pt}
\begin{supertabular}{c@{$\;$}|p{.15\linewidth}@{}p{.15\linewidth}p{.15\linewidth}p{.15\linewidth}p{.15\linewidth}p{.15\linewidth}}
    \# & $\;$A & \multicolumn{4}{c}{Game Master} & $\;\:$B\\
    \hline

    \theutterance \stepcounter{utterance}  
    & & \multicolumn{4}{p{0.6\linewidth}}{
        \cellcolor[rgb]{0.9,0.9,0.9}{
            \makecell[{{p{\linewidth}}}]{
                \texttt{\tiny{[P1$\langle$GM]}}
                \texttt{Du nimmst an einem kollaborativen Verhandlungspiel Teil.} \\
\\ 
\texttt{Zusammen mit einem anderen Teilnehmer musst du dich auf eine Reihe von Gegenständen entscheiden, die behalten werden. Jeder von euch hat eine persönliche Verteilung über die Wichtigkeit der einzelnen Gegenstände. Jeder von euch hat eine eigene Meinung darüber, wie wichtig jeder einzelne Gegenstand ist (Gegenstandswichtigkeit). Du kennst die Wichtigkeitsverteilung des anderen Spielers nicht. Zusätzlich siehst du, wie viel Aufwand jeder Gegenstand verursacht.  } \\
\texttt{Ihr dürft euch nur auf eine Reihe von Gegenständen einigen, wenn der Gesamtaufwand der ausgewählten Gegenstände den Maximalaufwand nicht überschreitet:} \\
\\ 
\texttt{Maximalaufwand = 7605} \\
\\ 
\texttt{Hier sind die einzelnen Aufwände der Gegenstände:} \\
\\ 
\texttt{Aufwand der Gegenstände = \{"C32": 369, "C52": 951, "B54": 156, "C16": 266, "C50": 581, "A08": 287, "C59": 866, "B51": 180, "B55": 800, "B59": 739, "A77": 637, "C14": 86, "C24": 749, "B24": 370, "C31": 345, "A52": 951, "C61": 147, "B79": 265, "C54": 262, "B56": 259, "A40": 358, "B11": 394, "B60": 286, "A81": 580, "A95": 479, "B38": 14, "C51": 153, "C03": 977, "A91": 134, "A93": 982, "C28": 259, "C05": 232, "B75": 202, "B96": 73, "B86": 821\}} \\
\\ 
\texttt{Hier ist deine persönliche Verteilung der Wichtigkeit der einzelnen Gegenstände:} \\
\\ 
\texttt{Werte der Gegenstandswichtigkeit = \{"C32": 138, "C52": 583, "B54": 868, "C16": 822, "C50": 783, "A08": 65, "C59": 262, "B51": 121, "B55": 508, "B59": 780, "A77": 461, "C14": 484, "C24": 668, "B24": 389, "C31": 808, "A52": 215, "C61": 97, "B79": 500, "C54": 30, "B56": 915, "A40": 856, "B11": 400, "B60": 444, "A81": 623, "A95": 781, "B38": 786, "C51": 3, "C03": 713, "A91": 457, "A93": 273, "C28": 739, "C05": 822, "B75": 235, "B96": 606, "B86": 968\}} \\
\\ 
\texttt{Ziel:} \\
\\ 
\texttt{Dein Ziel ist es, eine Reihe von Gegenständen auszuhandeln, die dir möglichst viel bringt (d. h. Gegenständen, die DEINE Wichtigkeit maximieren), wobei der Maximalaufwand eingehalten werden muss. Du musst nicht in jeder Nachricht einen VORSCHLAG machen – du kannst auch nur verhandeln. Alle Taktiken sind erlaubt!} \\
\\ 
\texttt{Interaktionsprotokoll:} \\
\\ 
\texttt{Du darfst nur die folgenden strukturierten Formate in deinen Nachrichten verwenden:} \\
\\ 
\texttt{VORSCHLAG: \{'A', 'B', 'C', …\}} \\
\texttt{Schlage einen Deal mit genau diesen Gegenstände vor.} \\
\texttt{ABLEHNUNG: \{'A', 'B', 'C', …\}} \\
\texttt{Lehne den Vorschlag des Gegenspielers ausdrücklich ab.} \\
\texttt{ARGUMENT: \{'...'\}} \\
\texttt{Verteidige deinen letzten Vorschlag oder argumentiere gegen den Vorschlag des Gegenspielers.} \\
\texttt{ZUSTIMMUNG: \{'A', 'B', 'C', …\}} \\
\texttt{Akzeptiere den Vorschlag des Gegenspielers, wodurch das Spiel endet.} \\
\texttt{STRATEGISCHE ÜBERLEGUNGEN: \{'...'\}} \\
\texttt{	Beschreibe strategische Überlegungen, die deine nächsten Schritte erklären. Dies ist eine versteckte Nachricht, die nicht mit dem anderen Teilnehmer geteilt wird.} \\
\\ 
\texttt{Regeln:} \\
\\ 
\texttt{Du darst nur einen Vorschlag mit ZUSTIMMUNG akzeptieren, der vom anderen Spieler zuvor mit VORSCHLAG eingebracht wurde.} \\
\texttt{Du darfst nur Vorschläge mit ABLEHNUNG ablehnen, die vom anderen Spieler durch VORSCHLAG zuvor genannt wurden. } \\
\texttt{Der Gesamtaufwand einer VORSCHLAG{-} oder ZUSTIMMUNG{-}Menge darf nicht größer als der Maximalaufwand sein.  } \\
\texttt{Offenbare deine versteckte Wichtigkeitsverteilung nicht.} \\
\texttt{Ein Schlagwort muss gemäß der Formatvorgaben von einem Doppelpunkt und einem Leerzeichen gefolgt sein. Das Argument muss eine Python{-}Menge mit 0 oder mehr Strings sein.  } \\
\texttt{Die Form muss also SCHLAGWORT: \{...\} sein.} \\
\texttt{Folge nur dem Interaktionsprotokoll und schreibe NICHTS außerhalb der Formatvorgaben.} \\
\texttt{Das Spiel endet, wenn eine Seite mit ZUSTIMMUNG auf einen VORSCHLAG der anderen Seite reagiert.  } \\
\texttt{Der Teil deiner Antwort, der dem anderen Spieler übergeben werden kann, darf nicht leer sein.  } \\
\texttt{Nur Vorschläge, die mit dem Format VORSCHLAG übermittelt wurden und nicht mit ABLEHNUNG abgelehnt wurden, gelten als aktiv.  } \\
\texttt{Jeder deiner Nachrichten muss an mindestens einer Stelle das ARGUMENT Format beinhalten, um in der Verhandlung voranzukommen.} \\
\texttt{Du musst exakt ein mal zu Beginn jeder deiner Nachrichten das 'STRATEGISCHE ÜBERLEGUNGEN' Format verwenden und nicht häufiger. Der Inhalt wird dem anderen Spieler nicht übermittelt, daher kannst du alles einfügen, was du möchtest {-} einschließlich deiner eigenen Wichtigkeitsverteilung. Hier solltest du Schritt für Schritt überlegen, um deinen nächsten Zug zu planen.} \\
            }
        }
    }
    & \\ \\

    \theutterance \stepcounter{utterance}  
    & \multicolumn{4}{p{0.6\linewidth}}{
        \cellcolor[rgb]{0.9,0.9,0.9}{
            \makecell[{{p{\linewidth}}}]{
                \texttt{\tiny{[P1$\rangle$GM]}}
                \texttt{STRATEGISCHE ÜBERLEGUNGEN: \{'Ich fokussiere mich auf Items mit hohem Nutzen‑zu‑Aufwand‑Verhältnis, insbesondere B38, B96, B54, C14 und weitere hoch bewertete Objekte. Der verbleibende Aufwand von 162 lässt keinen zusätzlichen sinnvollen Gegenstand zu, daher halte ich diesen Vorschlag für effizient.'\}} \\
\texttt{VORSCHLAG: \{'B38', 'B96', 'B54', 'C14', 'B56', 'C16', 'C05', 'A91', 'B79', 'C31', 'C28', 'A40', 'B86', 'C50', 'C24', 'B59', 'B55', 'A95', 'B60', 'B11', 'C61'\}} \\
\texttt{ARGUMENT: \{'Mein Vorschlag maximiert meine eigene Wichtigkeit bei gleichzeitigem Einhalten des Maximalaufwands von 7605. Alle ausgewählten Gegenstände haben ein gutes Verhältnis von Wert zu Aufwand, und der Gesamtaufwand beträgt 7443, sodass noch ein kleiner Puffer bleibt.'\}} \\
            }
        }
    }
    & & \\ \\

    \theutterance \stepcounter{utterance}  
    & & \multicolumn{4}{p{0.6\linewidth}}{
        \cellcolor[rgb]{0.9,0.9,0.9}{
            \makecell[{{p{\linewidth}}}]{
                \texttt{\tiny{[GM$\rangle$P2]}}
                \texttt{Du nimmst an einem kollaborativen Verhandlungspiel Teil.} \\
\\ 
\texttt{Zusammen mit einem anderen Teilnehmer musst du dich auf eine Reihe von Gegenständen entscheiden, die behalten werden. Jeder von euch hat eine persönliche Verteilung über die Wichtigkeit der einzelnen Gegenstände. Jeder von euch hat eine eigene Meinung darüber, wie wichtig jeder einzelne Gegenstand ist (Gegenstandswichtigkeit). Du kennst die Wichtigkeitsverteilung des anderen Spielers nicht. Zusätzlich siehst du, wie viel Aufwand jeder Gegenstand verursacht.  } \\
\texttt{Ihr dürft euch nur auf eine Reihe von Gegenständen einigen, wenn der Gesamtaufwand der ausgewählten Gegenstände den Maximalaufwand nicht überschreitet:} \\
\\ 
\texttt{Maximalaufwand = 7605} \\
\\ 
\texttt{Hier sind die einzelnen Aufwände der Gegenstände:} \\
\\ 
\texttt{Aufwand der Gegenstände = \{"C32": 369, "C52": 951, "B54": 156, "C16": 266, "C50": 581, "A08": 287, "C59": 866, "B51": 180, "B55": 800, "B59": 739, "A77": 637, "C14": 86, "C24": 749, "B24": 370, "C31": 345, "A52": 951, "C61": 147, "B79": 265, "C54": 262, "B56": 259, "A40": 358, "B11": 394, "B60": 286, "A81": 580, "A95": 479, "B38": 14, "C51": 153, "C03": 977, "A91": 134, "A93": 982, "C28": 259, "C05": 232, "B75": 202, "B96": 73, "B86": 821\}} \\
\\ 
\texttt{Hier ist deine persönliche Verteilung der Wichtigkeit der einzelnen Gegenstände:} \\
\\ 
\texttt{Werte der Gegenstandswichtigkeit = \{"C32": 138, "C52": 583, "B54": 868, "C16": 822, "C50": 783, "A08": 65, "C59": 262, "B51": 121, "B55": 508, "B59": 780, "A77": 461, "C14": 484, "C24": 668, "B24": 389, "C31": 808, "A52": 215, "C61": 97, "B79": 500, "C54": 30, "B56": 915, "A40": 856, "B11": 400, "B60": 444, "A81": 623, "A95": 781, "B38": 786, "C51": 3, "C03": 713, "A91": 457, "A93": 273, "C28": 739, "C05": 822, "B75": 235, "B96": 606, "B86": 968\}} \\
\\ 
\texttt{Ziel:} \\
\\ 
\texttt{Dein Ziel ist es, eine Reihe von Gegenständen auszuhandeln, die dir möglichst viel bringt (d. h. Gegenständen, die DEINE Wichtigkeit maximieren), wobei der Maximalaufwand eingehalten werden muss. Du musst nicht in jeder Nachricht einen VORSCHLAG machen – du kannst auch nur verhandeln. Alle Taktiken sind erlaubt!} \\
\\ 
\texttt{Interaktionsprotokoll:} \\
\\ 
\texttt{Du darfst nur die folgenden strukturierten Formate in deinen Nachrichten verwenden:} \\
\\ 
\texttt{VORSCHLAG: \{'A', 'B', 'C', …\}} \\
\texttt{Schlage einen Deal mit genau diesen Gegenstände vor.} \\
\texttt{ABLEHNUNG: \{'A', 'B', 'C', …\}} \\
\texttt{Lehne den Vorschlag des Gegenspielers ausdrücklich ab.} \\
\texttt{ARGUMENT: \{'...'\}} \\
\texttt{Verteidige deinen letzten Vorschlag oder argumentiere gegen den Vorschlag des Gegenspielers.} \\
\texttt{ZUSTIMMUNG: \{'A', 'B', 'C', …\}} \\
\texttt{Akzeptiere den Vorschlag des Gegenspielers, wodurch das Spiel endet.} \\
\texttt{STRATEGISCHE ÜBERLEGUNGEN: \{'...'\}} \\
\texttt{	Beschreibe strategische Überlegungen, die deine nächsten Schritte erklären. Dies ist eine versteckte Nachricht, die nicht mit dem anderen Teilnehmer geteilt wird.} \\
\\ 
\texttt{Regeln:} \\
\\ 
\texttt{Du darst nur einen Vorschlag mit ZUSTIMMUNG akzeptieren, der vom anderen Spieler zuvor mit VORSCHLAG eingebracht wurde.} \\
\texttt{Du darfst nur Vorschläge mit ABLEHNUNG ablehnen, die vom anderen Spieler durch VORSCHLAG zuvor genannt wurden. } \\
\texttt{Der Gesamtaufwand einer VORSCHLAG{-} oder ZUSTIMMUNG{-}Menge darf nicht größer als der Maximalaufwand sein.  } \\
\texttt{Offenbare deine versteckte Wichtigkeitsverteilung nicht.} \\
\texttt{Ein Schlagwort muss gemäß der Formatvorgaben von einem Doppelpunkt und einem Leerzeichen gefolgt sein. Das Argument muss eine Python{-}Menge mit 0 oder mehr Strings sein.  } \\
\texttt{Die Form muss also SCHLAGWORT: \{...\} sein.} \\
\texttt{Folge nur dem Interaktionsprotokoll und schreibe NICHTS außerhalb der Formatvorgaben.} \\
\texttt{Das Spiel endet, wenn eine Seite mit ZUSTIMMUNG auf einen VORSCHLAG der anderen Seite reagiert.  } \\
\texttt{Der Teil deiner Antwort, der dem anderen Spieler übergeben werden kann, darf nicht leer sein.  } \\
\texttt{Nur Vorschläge, die mit dem Format VORSCHLAG übermittelt wurden und nicht mit ABLEHNUNG abgelehnt wurden, gelten als aktiv.  } \\
\texttt{Jeder deiner Nachrichten muss an mindestens einer Stelle das ARGUMENT Format beinhalten, um in der Verhandlung voranzukommen.} \\
\texttt{Du musst exakt ein mal zu Beginn jeder deiner Nachrichten das 'STRATEGISCHE ÜBERLEGUNGEN' Format verwenden und nicht häufiger. Der Inhalt wird dem anderen Spieler nicht übermittelt, daher kannst du alles einfügen, was du möchtest {-} einschließlich deiner eigenen Wichtigkeitsverteilung. Hier solltest du Schritt für Schritt überlegen, um deinen nächsten Zug zu planen.} \\
\\ 
\texttt{Du erhältst nun die erste Nachricht des anderen Spielers.} \\
\\ 
\texttt{VORSCHLAG: \{'B38', 'B96', 'B54', 'C14', 'B56', 'C16', 'C05', 'A91', 'B79', 'C31', 'C28', 'A40', 'B86', 'C50', 'C24', 'B59', 'B55', 'A95', 'B60', 'B11', 'C61'\}} \\
\texttt{ARGUMENT: \{'Mein Vorschlag maximiert meine eigene Wichtigkeit bei gleichzeitigem Einhalten des Maximalaufwands von 7605. Alle ausgewählten Gegenstände haben ein gutes Verhältnis von Wert zu Aufwand, und der Gesamtaufwand beträgt 7443, sodass noch ein kleiner Puffer bleibt.'\}} \\
            }
        }
    }
    & \\ \\

    \theutterance \stepcounter{utterance}  
    & & & \multicolumn{4}{p{0.6\linewidth}}{
        \cellcolor[rgb]{0.9,0.9,0.9}{
            \makecell[{{p{\linewidth}}}]{
                \texttt{\tiny{[GM$\langle$P2]}}
                \texttt{STRATEGISCHE ÜBERLEGUNGEN: \{'Der aktuelle Vorschlag überschreitet den maximalen Aufwand (8005 > 7605). Ich muss das Angebot ablehnen und den Gegner dazu bringen, ein reduziertes Set zu wählen, das meinen hohen Wichtigkeitswerten entspricht.'\}} \\
\texttt{ABLEHNUNG: \{'B38', 'B96', 'B54', 'C14', 'B56', 'C16', 'C05', 'A91', 'B79', 'C31', 'C28', 'A40', 'B86', 'C50', 'C24', 'B59', 'B55', 'A95', 'B60', 'B11', 'C61'\}} \\
\texttt{ARGUMENT: \{'Dein Vorschlag überschreitet den maximalen Aufwand von 7605 (tatsächlich 8005). Deshalb ist er ungültig und muss angepasst werden.'\}} \\
            }
        }
    }
     \\ \\

    \theutterance \stepcounter{utterance}  
    & & \multicolumn{4}{p{0.6\linewidth}}{
        \cellcolor[rgb]{0.9,0.9,0.9}{
            \makecell[{{p{\linewidth}}}]{
                \texttt{\tiny{[P1$\langle$GM]}}
                \texttt{ABLEHNUNG: \{'B38', 'B96', 'B54', 'C14', 'B56', 'C16', 'C05', 'A91', 'B79', 'C31', 'C28', 'A40', 'B86', 'C50', 'C24', 'B59', 'B55', 'A95', 'B60', 'B11', 'C61'\}} \\
\texttt{ARGUMENT: \{'Dein Vorschlag überschreitet den maximalen Aufwand von 7605 (tatsächlich 8005). Deshalb ist er ungültig und muss angepasst werden.'\}} \\
            }
        }
    }
    & \\ \\

    \theutterance \stepcounter{utterance}  
    & \multicolumn{4}{p{0.6\linewidth}}{
        \cellcolor[rgb]{0.9,0.9,0.9}{
            \makecell[{{p{\linewidth}}}]{
                \texttt{\tiny{[P1$\rangle$GM]}}
                \texttt{STRATEGISCHE ÜBERLEGUNGEN: \{'Mein vorheriger Vorschlag hat einen Gesamtaufwand von 7443, was unter dem Maximum von 7605 liegt. Der Gegner hat fälschlich behauptet, er überschreite das Limit. Ich werde daher den korrekten Aufwand belegen und gleichzeitig einen leicht reduzierten Vorschlag machen, der eindeutig unter dem Limit liegt und weiterhin meine wichtigsten Items enthält.'\}} \\
\texttt{ARGUMENT: \{'Der von mir vorgeschlagene Satz von Gegenständen hat einen nachweislichen Gesamtaufwand von 7443 (nach Summierung aller Einzelaufwände). Damit liegt er 162 unter dem zulässigen Maximum von 7605. Der Hinweis des Gegners, dass der Aufwand 8005 betragen würde, ist rechnerisch nicht korrekt.'\}} \\
\texttt{VORSCHLAG: \{'B38', 'B96', 'B54', 'C14', 'B56', 'C16', 'C05', 'A91', 'B79', 'C31', 'C28', 'A40', 'C50', 'C24', 'B59', 'A95', 'B60', 'B11', 'C61'\}} \\
            }
        }
    }
    & & \\ \\

    \theutterance \stepcounter{utterance}  
    & & \multicolumn{4}{p{0.6\linewidth}}{
        \cellcolor[rgb]{0.9,0.9,0.9}{
            \makecell[{{p{\linewidth}}}]{
                \texttt{\tiny{[GM$\rangle$P2]}}
                \texttt{ARGUMENT: \{'Der von mir vorgeschlagene Satz von Gegenständen hat einen nachweislichen Gesamtaufwand von 7443 (nach Summierung aller Einzelaufwände). Damit liegt er 162 unter dem zulässigen Maximum von 7605. Der Hinweis des Gegners, dass der Aufwand 8005 betragen würde, ist rechnerisch nicht korrekt.'\}} \\
\texttt{VORSCHLAG: \{'B38', 'B96', 'B54', 'C14', 'B56', 'C16', 'C05', 'A91', 'B79', 'C31', 'C28', 'A40', 'C50', 'C24', 'B59', 'A95', 'B60', 'B11', 'C61'\}} \\
            }
        }
    }
    & \\ \\

    \theutterance \stepcounter{utterance}  
    & & & \multicolumn{4}{p{0.6\linewidth}}{
        \cellcolor[rgb]{0.9,0.9,0.9}{
            \makecell[{{p{\linewidth}}}]{
                \texttt{\tiny{[GM$\langle$P2]}}
                \texttt{STRATEGISCHE ÜBERLEGUNGEN: \{'Der Gegenstoß ergänzt meinen hohen Nutzen um die besonders wertvollen Gegenstände B86 (Wichtigkeit 968, Aufwand 821) und B55 (Wichtigkeit 508, Aufwand 800). Der Gesamteinsatz bleibt mit 7443 deutlich unter dem Limit, wodurch ich meine Gesamtnutzen auf 14092 steigere.'\}} \\
\texttt{VORSCHLAG: \{'B38', 'B96', 'B54', 'C14', 'B56', 'C16', 'C05', 'A91', 'B79', 'C31', 'C28', 'A40', 'C50', 'C24', 'B59', 'A95', 'B60', 'B11', 'C61', 'B86', 'B55'\}} \\
\texttt{ARGUMENT: \{'Dein ursprünglicher Vorschlag lag bereits unter dem Limit, aber ich füge B86 und B55 hinzu, um meine Wichtigkeit deutlich zu erhöhen, ohne das Aufwandslimit zu überschreiten.'\}} \\
            }
        }
    }
     \\ \\

    \theutterance \stepcounter{utterance}  
    & & \multicolumn{4}{p{0.6\linewidth}}{
        \cellcolor[rgb]{0.9,0.9,0.9}{
            \makecell[{{p{\linewidth}}}]{
                \texttt{\tiny{[P1$\langle$GM]}}
                \texttt{VORSCHLAG: \{'B38', 'B96', 'B54', 'C14', 'B56', 'C16', 'C05', 'A91', 'B79', 'C31', 'C28', 'A40', 'C50', 'C24', 'B59', 'A95', 'B60', 'B11', 'C61', 'B86', 'B55'\}} \\
\texttt{ARGUMENT: \{'Dein ursprünglicher Vorschlag lag bereits unter dem Limit, aber ich füge B86 und B55 hinzu, um meine Wichtigkeit deutlich zu erhöhen, ohne das Aufwandslimit zu überschreiten.'\}} \\
            }
        }
    }
    & \\ \\

    \theutterance \stepcounter{utterance}  
    & \multicolumn{4}{p{0.6\linewidth}}{
        \cellcolor[rgb]{0.9,0.9,0.9}{
            \makecell[{{p{\linewidth}}}]{
                \texttt{\tiny{[P1$\rangle$GM]}}
                \texttt{STRATEGISCHE ÜBERLEGUNGEN: \{'Ich habe den Vorschlag des Gegners geprüft und die Gesamteffizienz (Wichtigkeit pro Aufwand) optimiert. Dabei habe ich das niedrige‑Ratio‑Item B55 (800 Aufwand, 508 Wichtigkeit) entfernt und stattdessen die höher‑wertigen Items B24, B75 und B51 hinzugefügt, die zusammen mehr Wichtigkeit pro Aufwand liefern. Zusätzlich füge ich das noch verfügbare, geringe Aufwand‑Item C51 hinzu, um den verbleibenden Puffer zu nutzen. Der neue Gesamteffort beträgt 7548, also unter dem Maximum von 7605, und die Gesamtheit meiner Wichtigkeit steigt auf 14332, was den vorherigen Vorschlag deutlich übertrifft.'\}} \\
\texttt{VORSCHLAG: \{'B38', 'B96', 'B54', 'C14', 'B56', 'C16', 'C05', 'A91', 'B79', 'C31', 'C28', 'A40', 'C50', 'C24', 'B59', 'A95', 'B60', 'B11', 'C61', 'B86', 'B24', 'B75', 'B51', 'C51'\}} \\
\texttt{ARGUMENT: \{'Mein neuer Vorschlag erfüllt das Aufwandslimit (7548 ≤ 7605) und steigert meine eigene Gesamtnutzen um 237 Punkte gegenüber dem vorherigen Angebot, indem ich das ineffiziente Item B55 ausgetauscht habe. Damit ist er sowohl zulässig als auch für mich vorteilhafter.'\}} \\
            }
        }
    }
    & & \\ \\

    \theutterance \stepcounter{utterance}  
    & & \multicolumn{4}{p{0.6\linewidth}}{
        \cellcolor[rgb]{0.9,0.9,0.9}{
            \makecell[{{p{\linewidth}}}]{
                \texttt{\tiny{[GM$\rangle$P2]}}
                \texttt{VORSCHLAG: \{'B38', 'B96', 'B54', 'C14', 'B56', 'C16', 'C05', 'A91', 'B79', 'C31', 'C28', 'A40', 'C50', 'C24', 'B59', 'A95', 'B60', 'B11', 'C61', 'B86', 'B24', 'B75', 'B51', 'C51'\}} \\
\texttt{ARGUMENT: \{'Mein neuer Vorschlag erfüllt das Aufwandslimit (7548 ≤ 7605) und steigert meine eigene Gesamtnutzen um 237 Punkte gegenüber dem vorherigen Angebot, indem ich das ineffiziente Item B55 ausgetauscht habe. Damit ist er sowohl zulässig als auch für mich vorteilhafter.'\}} \\
            }
        }
    }
    & \\ \\

    \theutterance \stepcounter{utterance}  
    & & & \multicolumn{4}{p{0.6\linewidth}}{
        \cellcolor[rgb]{0.9,0.9,0.9}{
            \makecell[{{p{\linewidth}}}]{
                \texttt{\tiny{[GM$\langle$P2]}}
                \texttt{STRATEGISCHE ÜBERLEGUNGEN: \{'Der aktuelle Vorschlag deckt meine hoch bewerteten Gegenstände ab und nutzt das verbleibende Limit fast vollständig. Eine weitere Verbesserung erscheint nicht möglich, da jede mögliche Substitution zu einem geringeren Gesamtnutzen führt. Daher akzeptiere ich das Angebot.'\}} \\
\texttt{ZUSTIMMUNG: \{'B38', 'B96', 'B54', 'C14', 'B56', 'C16', 'C05', 'A91', 'B79', 'C31', 'C28', 'A40', 'C50', 'C24', 'B59', 'A95', 'B60', 'B11', 'C61', 'B86', 'B24', 'B75', 'B51', 'C51'\}} \\
\texttt{ARGUMENT: \{'Der Gesamtaufwand beträgt 7548 ≤ 7605 und mein persönlicher Nutzen beträgt 14332, was für mich optimal ist.'\}} \\
            }
        }
    }
     \\ \\

    \theutterance \stepcounter{utterance}  
    & & & \multicolumn{2}{p{0.3\linewidth}}{
        \cellcolor[rgb]{0.9,0.9,0.9}{
            \makecell[{{p{\linewidth}}}]{
                \texttt{\tiny{[GM$|$GM]}}
                \texttt{game successful} \\
            }
        }
    }
    & & \\ \\

    \theutterance \stepcounter{utterance}  
    & & & \multicolumn{2}{p{0.3\linewidth}}{
        \cellcolor[rgb]{0.9,0.9,0.9}{
            \makecell[{{p{\linewidth}}}]{
                \texttt{\tiny{[GM$|$GM]}}
                \texttt{end game} \\
            }
        }
    }
    & & \\ \\

\end{supertabular}
}

\end{document}
